\scanpage{35}
% archive sheet number: 34
$\operatorname{Aff}_A(G)$ opère \deleted{à gauche et à droite} sur $U_A(G)$ ; \unsure{les formules sont|la formule est} \deleted{\unsure{de façon|de la façon|transposée}} \unsure{analogues|analogue}

\[
\langle g,f\mathbin{\lrcorner}X\rangle=\langle gf,X\rangle
\qquad\text{en particulier}\quad\langle f,X\rangle=\varepsilon(f\mathbin{\lrcorner}X)
\]

Supposons maintenant que $A\subset B$, alors

\[
\operatorname{Aff}_B(G)=\operatorname{Aff}_A(G)\otimes_A B\supset\operatorname{Aff}_A(G)
\]

[on a une \uline{inclusion}, car $\operatorname{Aff}_A(G)$ plat sur $A$]

\[
U_B(G)=U_A(G)\otimes_A B\supset U_A(G)
\]

[raison analogue ; la première égalité est due au fait que $G$ est $A$-simple]

\uline{Proposition 1}\quad Soit $X\in U_B(G)$. Conditions équivalentes

\[
\begin{aligned}
\text{(i)}\quad &X\in U_A(G)\\
\text{(ii)}\quad &X*\operatorname{Aff}_A(G)\subset\operatorname{Aff}_A(G)
  \qquad\text{(ii bis)}\quad\operatorname{Aff}_A(G)*X\subset\operatorname{Aff}_A(G)\\
\text{(iii)}\quad &\langle X,\operatorname{Aff}_A(G)\rangle\subset A
\end{aligned}
\]

En effet,

\diagram{\begin{tikzcd}[ampersand replacement=\&,column sep=small,row sep=tiny]
\& \text{(ii bis)}\arrow[dr,Rightarrow] \& {}\\
\text{(i)}\arrow[ur,Rightarrow]\arrow[r,Rightarrow] \& \text{(ii)}\arrow[r,Rightarrow] \& \text{(iii)}
\end{tikzcd}}

triviaux.

(iii) implique que \deleted{$X\unsure{\in|\mathbin{\lrcorner}}\operatorname{Aff}$} \deleted{la forme} \inserted{$X$ \unsure{définit|est}} une forme $A$-linéaire sur $\operatorname{Aff}_A(G)$ nulle sur une puissance de l'idéal d'augmentation, donc un élément $X_{\unsure{0|o}}$ de $U_A(G)$, qui n'est évidemment autre que $X$
